% =====================================================================
% RESUMEN
% =====================================================================
% Todas las cifras de este resumen estan verificadas contra resultados.py por
% scripts/figuras/verificar_memoria.py. Si se reescribe la redaccion, no se
% deben tocar los numeros sin actualizar antes esa fuente.
%
% Solo en castellano: el autor confirma que la normativa no exige version
% en euskera ni en ingles. La inglesa que hubo aqui queda en git.
%
% ESTILO: sin negritas en el texto corrido y sin incisos entre rayas.
% =====================================================================

Los sistemas de detección de intrusiones basados en aprendizaje automático
alcanzan resultados casi perfectos sobre los conjuntos de datos públicos de
referencia. Este trabajo parte de una pregunta incómoda sobre esas cifras: qué
miden realmente. Para responderla se construye un \emph{pipeline} completo que
convierte capturas de red en bruto en ejemplos etiquetados, en total 17.203.272
conexiones TCP repartidas en seis días del conjunto CSE-CIC-IDS2018, y se
comparan tres representaciones del tráfico: la distribución de bytes de la carga
útil, las características de flujo y una vista conductual basada en la ráfaga de
conexiones por origen. Se procesan además tres días de CIC-IDS2017 con los
ataques análogos, que sirven para la validación entre conjuntos de datos.

La primera aportación es metodológica. La \emph{ground-truth} no se hereda de la
documentación del conjunto de datos, sino que se reconstruye sobre las propias
capturas, y al hacerlo se detecta que la documentación oficial es incorrecta en
tres de los siete días analizados. Se documenta además un error que este trabajo
cometió dos veces: el tráfico benigno de referencia estaba contaminado por
conexiones del propio atacante, lo que producía resultados plausibles y sin
sentido.

El resultado central es negativo y reordena el trabajo. Medidas las tres
representaciones con la misma métrica y el mismo protocolo, ninguna baja de
$0.9186$ y el análisis de contenido no baja de $0.9734$, de modo que la
evaluación dentro de un mismo día está saturada y no discrimina entre ellas. Se
propone una explicación, que un único host atacante por día permita memorizar su
identidad, y se contrasta y refuta mediante un experimento de exclusión de
atacantes con control positivo. Lo que las representaciones aprenden no es el
host, sino la herramienta con que se generó el ataque.

La discriminación aparece al perturbar las condiciones. Camuflar los primeros
bytes de cada flujo hunde el análisis de contenido de un \emph{recall} de
$1.0000$ a $0.0025$, salvo en el régimen en claro, donde la firma es
contenido genuino. Entre días distintos, lo que falla no es la representación
sino el umbral: el AUC se mantiene por encima de $0.9724$ y bastan 25 flujos
etiquetados del dominio nuevo para recuperar el rendimiento sin reentrenar. Y se
muestra que la validación entre días no detecta la fuga de identificadores de
red, sino que la premia: solo el cambio a otro conjunto de datos la rompe, y al
hacerlo revela que dos de los tres regímenes transfieren por la huella de la
herramienta y no por el ataque.

Finalmente, se contrastan dos arquitecturas del estado del arte, un
\emph{transformer} sobre bytes y una red de atención cruzada, a presupuesto de
parámetros comparable. Ninguna mejora los resultados, y cuestan entre 2 y 49
veces más. La conclusión del trabajo es que, sobre estos datos, el cuello de
botella no está en la capacidad del modelo sino en la calidad de la etiqueta y
en el protocolo con que se mide.
